Physical AI systems such as robots, drones, and autonomous vehicles run their perception stacks on GPUs. For these systems, the dominant cost is increasingly not arithmetic, but data movement: the time and energy spent transferring bytes between memory and compute resources. Memory-centric hardware including processing-in-memory (PiM), in/near-storage processing, and near-sensor computation aims to address this cost. However, hardware cannot be optimally designed for workloads that have not been measured.

This thesis presents the first per-kernel, per-data-structure memory characterization of a \emph{production} visual-SLAM system: cuVSLAM, the CUDA-accelerated simultaneous localization and mapping library deployed in NVIDIA's Isaac robotics stack.

The thesis makes three primary contributions:
\begin{enumerate}
\item \textbf{Methodology:} It introduces GPU-DAMOV, a re-derivation of the DAMOV data-movement methodology for GPUs. In this approach, the last-to-first miss ratio is redefined for a two-level cache hierarchy. Latency-boundedness is conditioned on occupancy, and memory-divergence (coalescing) is introduced as a first-class axis with no CPU analogue.

\item \textbf{Empirical Analysis:} The method is applied to cuVSLAM across 27 sequences from four public datasets (KITTI, EuRoC, TUM RGB-D, TUM-VI), with zero failed runs. This yields an eight-class bottleneck taxonomy (G0–G7) that is stable across inputs (91\% cross-sequence class consistency), independently validated by two unsupervised clustering algorithms (k-means and Ward hierarchical; purity $\approx$0.68), and 80\% reproducible across two GPU microarchitectures.

\item \textbf{Trust and Validation:} A novel allocation-journaling instrumentation (TaggedAllocator), combined with binary-instrumentation address traces, attributes every byte of DRAM traffic to a named data structure. All 274 allocations are resolved, and 48 of 49 kernels carry a unanimous data-structure tag across all sequences.
\end{enumerate}

Additionally, the thesis demonstrates that the profiled system's trajectory accuracy reproduces NVIDIA’s published results across a 141-configuration matrix. Profiling itself is shown to be accuracy-neutral on deterministic modes: the profiled trajectories are bit-identical to unprofiled runs across a 192-configuration campaign.


The characterization resolves cuVSLAM's memory behaviour into a three-way
split with direct architectural consequences. The image front-end is
\emph{cache-immune streaming}: its measured reuse-distance profile is flat
from 64\,KiB to 48\,MiB, so no realizable cache helps, and the measured
1.68\,MB-per-frame sensor upload argues for near-sensor consumption. The
bundle-adjustment solver is hot-persistent and largely compute-efficient: its
DRAM traffic is 97\% one named, pre-sized linear-system structure. The
loop-closure back-end is a scattered gather (23--30 sectors per warp access)
over a keyframe database whose GPU-resident state is a fixed 6.7\,MB buffer
while the session-scale store grows host-side (the in-storage-processing
case); and 92.6\% of the scan kernel's DRAM volume is register-spill traffic,
a compiler artifact invisible to counter-only studies. In total, 60--72\% of
GPU time carries PiM affinity under a sensitivity-tested classification. A
measured whole-run energy baseline (34.7\,J) and host-side I/O
characterization (66\,MB storage ingestion; 708\,MB peak host memory) complete
the evidence base for the follow-on architecture study, for which
simulator-configuration groundwork is laid.

Beyond its findings, the thesis documents two self-corrections in which
independent measurement methods initially disagreed and were reconciled (one
reversing a claim made earlier in the project), and argues that this
disagree-then-agree arc, with all artifacts on record, is precisely what
distinguishes measurement science from advocacy.
